\section{Setup} 
\label{language_modeling}
Pretrained language understanding models are central tasks in natural language processing and language understanding. There are several formulations of language modeling. In this work we focus on GPT-2 \cite{Radford2019GPT2}, a left-to-right generative transformer based language model, and BERT \cite{devlin2018bert}, a bi-directional transformer model based on language model masking. We explain our configurations for these models in the following section and refer to the original papers for more details.

\subsection{Training Dataset}

To collect a large diverse training set with longterm dependencies we aggregate several of the largest language modeling datasets. We create an aggregate dataset consisting of Wikipedia \citep{devlin2018bert}, CC-Stories \citep{ccstories}, RealNews \citep{grover}, and OpenWebtext \cite{Radford2019GPT2}. To avoid training set leakage into our downstream tasks we remove the Wikipedia articles present in the WikiText103 test set \citep{wikitext}. We  also remove unnecessary newlines from the CC-Stories corpus introduced by preprocessing artifacts. For BERT models we include BooksCorpus \cite{BooksCorpus} in the training dataset, however, this dataset is excluded for GPT-2 trainings as it overlaps with LAMBADA task.

We combined all the datasets and then filtered out all the documents with content length less than 128 tokens from the aggregated dataset. Since similar content might appear multiple times in the aggregated datasets, we used locality-sensitive hashing (LSH) to deduplicate content with a jaccard similarity greater than 0.7. The resulting aggregate corpus contains 174 GB of deduplicated text.

\subsection{Training Optimization and Hyperparameters}
To train our models efficiently we utilize mixed precision training with dynamic loss scaling to take advantage of the V100's Tensor Cores \citep{MPTraining, AutoLossScale}. We start by initializing our weights $W$ with a simple normal distribution $W \sim \mathcal{N}(0, 0.02)$. We then scale weights immediately before residual layers by $\frac{1}{\sqrt{2N}}$ where N is the number of transformer layers comprised of self attention and MLP blocks. For our optimizer we utilize  Adam  \citep{Adam} with weight decay \citep{loshchilov2018decoupled} $\lambda = 0.01$. 
Additionally, we use global gradient norm clipping of 1.0 to improve the stability of training large models. In all cases, a dropout of 0.1 is used.
 Lastly, to better manage our memory footprint we utilize activation checkpointing \citep{activation_checkpointing} after every transformer layer.
 
%\subsubsection{GPT-2}

For GPT-2 models, all training is performed with sequences of 1024 subword units at a batch size of 512 for 300k iterations. Our learning rate of 1.5e-4 utilizes a warmup period of 3k iterations before following a single cycle cosine decay over the remaining 297k iterations. We stop the decay at a minimum learning rate of 1e-5.

%\subsubsection{BERT}
For BERT models, we largely follow the training process described in \cite{ALBERT2019}. We use the original BERT dictionary with vocab size of 30,522. In addition, we replace the next sentence prediction head with sentence order prediction as suggested by \cite{ALBERT2019} and use whole word n-gram masking of \cite{SpanBERT2019}. For all cases, we set the batch size to 1024 and use a learning rate of 1.0e-4 warmed up over 10,000 iterations and decayed linearly over 2 million iterations. Other training parameters are kept the same as \cite{devlin2018bert}.



